% ==STYLE-BEGIN==
\documentclass[UTF8,11pt]{ctexart}
\usepackage{amsmath,amssymb,amsthm}
\usepackage[a4paper,margin=2.4cm,top=2.3cm,bottom=2.3cm]{geometry}
\usepackage{fancyhdr}
\usepackage{enumitem}
\usepackage{hyperref}
\hypersetup{hidelinks}

% ---- 版式 ----
\linespread{1.15}
\setlength{\parskip}{0.35em}
\setlength{\parindent}{1.6em}
\setlist{nosep,leftmargin=2em}

% ---- 页眉 ----
\pagestyle{fancy}
\fancyhf{}
\fancyhead[L]{\small\coursename}
\fancyhead[R]{\small\thepage}
\renewcommand{\headrulewidth}{0.4pt}

% ---- 定理环境：中文用正体，标题加粗，按小节共用一条流水号 ----
\newtheoremstyle{cnthm}{0.7em}{0.7em}{\normalfont}{}{\bfseries}{}{0.5em}{}
\theoremstyle{cnthm}
\newtheorem{thm}{定理}[section]
\newtheorem{lem}[thm]{引理}
\newtheorem{prop}[thm]{命题}
\newtheorem{cor}[thm]{推论}
\newtheorem{defn}[thm]{定义}
\newtheorem{rem}[thm]{注}
\newtheorem{exa}[thm]{例}
\newtheorem{conj}[thm]{猜想}
\renewcommand{\proofname}{\bfseries 证明}
\providecommand{\coursename}{}
% ==STYLE-END==
\renewcommand{\coursename}{解析几何\,·\,0924}
\title{几何与代数 0924}
\date{}
\begin{document}
\maketitle

\section{\ 例：$\vec d=x\vec a+y\vec b+z\vec c$ 的系数}

$\vec a,\vec b,\vec c$ 不共面，
$\vec d=x\vec a+y\vec b+z\vec c$，求 $x,y,z$。

若已知 $\vec a=(a_1,a_2,a_3)$，
$\vec b=(b_1,b_2,b_3)$，$\vec c=(c_1,c_2,c_3)$，
$\vec d=(d_1,d_2,d_3)$，上式等价于
\[
\begin{pmatrix}
d_1\\d_2\\d_3
\end{pmatrix}
=x\begin{pmatrix}
a_1\\a_2\\a_3
\end{pmatrix}
+y\begin{pmatrix}
b_1\\b_2\\b_3
\end{pmatrix}
+z\begin{pmatrix}
c_1\\c_2\\c_3
\end{pmatrix}
\]
\[
\Longleftrightarrow
\begin{cases}
d_1=xa_1+yb_1+zc_1,\\
d_2=xa_2+yb_2+zc_2,\\
d_3=xa_3+yb_3+zc_3 .
\end{cases}
\]
用之前求解三元一次方程组的方法即可写出 $x,y,z$ 的解，
此为代数解法；几何运用混合积的想法详见上节笔记。

\section{\ 平行坐标系}

在空间中，任取三个线性无关的向量组
$\vec a_1,\vec a_2,\vec a_3$，则称
$[O;\vec a_1,\vec a_2,\vec a_3]$ 是空间中的\textbf{平行坐标系}，
称 $\vec a_1,\vec a_2,\vec a_3$ 是\textbf{基向量}。此时
$\forall\vec b=x\vec a_1+y\vec a_2+z\vec a_3$，称
$(x,y,z)$ 是 $\vec b$ 在
$[O;\vec a_1,\vec a_2,\vec a_3]$ 下的坐标。

若 $\vec b=(x_1,y_1,z_1)$，$\vec c=(x_2,y_2,z_2)$，则
\[
\vec b+\vec c=(x_1+x_2,\ y_1+y_2,\ z_1+z_2),
\qquad
\lambda\vec b=(\lambda x_1,\ \lambda y_1,\ \lambda z_1).
\]
对于 $\vec b\cdot\vec c$，则有
\[
\begin{aligned}
\vec b\cdot\vec c
&=(x_1\vec a_1+y_1\vec a_2+z_1\vec a_3)
\cdot(x_2\vec a_1+y_2\vec a_2+z_2\vec a_3)\\
&=(x_1,y_1,z_1)
\begin{pmatrix}
\vec a_1\cdot\vec a_1 & \vec a_1\cdot\vec a_2 & \vec a_1\cdot\vec a_3\\
\vec a_2\cdot\vec a_1 & \vec a_2\cdot\vec a_2 & \vec a_2\cdot\vec a_3\\
\vec a_3\cdot\vec a_1 & \vec a_3\cdot\vec a_2 & \vec a_3\cdot\vec a_3
\end{pmatrix}
\begin{pmatrix}
x_2\\y_2\\z_2
\end{pmatrix}.
\end{aligned}
\]
$\forall i,j$，$g_{ij}=\vec a_i\cdot\vec a_j$，
我们称 $(g_{ij})$ 是\textbf{度量矩阵}。

\section{\ 对偶坐标系}

令 $\vec a_1'$ 是与 $\vec a_2,\vec a_3$ 垂直，且
$\vec a_1\cdot\vec a_1'=1$ 的向量；同理，$\vec a_2'$ 是与
$\vec a_1,\vec a_3$ 垂直，且 $\vec a_2\cdot\vec a_2'=1$；
$\vec a_3'$ 是与 $\vec a_1,\vec a_2$ 垂直，且
$\vec a_3\cdot\vec a_3'=1$。

设 $\vec a_1'=\lambda\vec a_2\times\vec a_3$，由
$\vec a_1\cdot\vec a_1'=1$ 可知
\[
\lambda=\frac{1}{\vec a_1\cdot\vec a_2\times\vec a_3},
\qquad
\vec a_1'=\frac{\vec a_2\times\vec a_3}
{\vec a_1\cdot\vec a_2\times\vec a_3},
\]
同理可知
\[
\vec a_2'=\frac{\vec a_3\times\vec a_1}
{\vec a_1\cdot\vec a_2\times\vec a_3},\qquad
\vec a_3'=\frac{\vec a_1\times\vec a_2}
{\vec a_1\cdot\vec a_2\times\vec a_3} .
\]

\textbf{性质 1：}这三个向量线性无关。由于
\[
(\vec a_2\times\vec a_3)\cdot
\bigl((\vec a_3\times\vec a_1)\times(\vec a_1\times\vec a_2)\bigr)
=(\vec a_1\cdot\vec a_2\times\vec a_3)^2\neq0,
\]
因此 $\vec a_1',\vec a_2',\vec a_3'$ 可作为一个新的坐标系的基向量。
称 $[O;\vec a_1',\vec a_2',\vec a_3']$ 是
$[O;\vec a_1,\vec a_2,\vec a_3]$ 的\textbf{对偶坐标系}。

\textbf{性质 2：}若 $[O;\vec a_1',\vec a_2',\vec a_3']$ 是
$[O;\vec a_1,\vec a_2,\vec a_3]$ 的对偶坐标系，则
$[O;\vec a_1',\vec a_2',\vec a_3']$ 的对偶坐标系是
$[O;\vec a_1,\vec a_2,\vec a_3]$。

直角坐标系 $[O;\vec e_1,\vec e_2,\vec e_3]$ 的对偶是它自己，
无论左右手系。

若 $\vec b$ 在 $[O;\vec a_1,\vec a_2,\vec a_3]$ 的坐标是
$(x,y,z)$，在 $[O;\vec a_1',\vec a_2',\vec a_3']$ 的坐标是
$(x',y',z')$，那么
\[
\vec b\cdot\vec b'=|\vec b|^2=xx'+yy'+zz' .
\]
若 $\vec c$ 在 $[O;\vec a_1,\vec a_2,\vec a_3]$ 中的坐标为
$(x_1,y_1,z_1)$，那么
$\vec b\cdot\vec c=xx_1+yy_1+zz_1$。

\subsubsection*{e.g.}

$\vec b$ 在 $[O;\vec a_1,\vec a_2,\vec a_3]$ 的坐标为
$(x,y,z)$，即
$\vec b=x\vec a_1+y\vec a_2+z\vec a_3$，不难看出
\[
\vec b\cdot\vec a_1'=x,\qquad
\vec b\cdot\vec a_2'=y,\qquad
\vec b\cdot\vec a_3'=z .
\]

\end{document}

\end{document}
